\documentclass[../../main.tex]{subfiles}
\graphicspath{{../../image/}{../../figure/}}
\begin{document}

\section{Liu 激波容许性准则}\label{section:ch08-liu容许性准则}

Liu 激波容许性检验比 Lax $E$-准则更精细，并力图刻画激波的内部稳定性。按其设计，它仅在由激波曲线连接状态的激波语境下才有意义。因此，它对一般系统的适用性限于中等强度激波。尽管如此，在特殊系统中它也适用于强激波。

对给定状态 $U_-$，设 $\lambda_i(U_-)$ 是 $\mathrm{D}F(U_-)$ 的单特征值，从而由\cref{theorem:ch08-激波曲线}，过 $U_-$ 的 $i$-激波曲线 $W_i(\tau;U_-)$ 有良好定义。一个在左侧连接 $U_-$、在右侧连接 $U_+=W_i(\tau_+;U_-)$、速度为 $s$ 的 $i$-激波满足 \emph{Liu $E$-条件}，如果
\begin{align}
s=s_i(\tau_+;U_-)\leqslant s_i(\tau;U_-),\qquad\text{对所有介于 }0\text{ 与 }\tau_+\text{ 之间的 }\tau.
\label{align:ch08-liu}
\end{align}

上述容许性准则的正当性将事后建立：既通过它与其他受物理驱动的激波容许性准则的联系，也通过它在构造 Riemann 问题稳定解（见\cref{chapter:ch09-容许波扇与riemann问题}）中的作用。

由于 $U_-$ 与 $U_+$ 由一个 $i$-激波相连，$U_-$ 也必须位于从 $U_+$ 发出的 $i$-激波曲线上，记作 $U_-=W_i(\tau_-;U_+)$。只要沿上述激波曲线\eqref{align:ch08-符号} 成立，\eqref{align:ch08-liu} 等价于对偶命题
\begin{align}
s=s_i(\tau_-;U_+)\geqslant s_i(\tau;U_+),\qquad\text{对所有介于 }0\text{ 与 }\tau_-\text{ 之间的 }\tau.
\label{align:ch08-liu对偶}
\end{align}

我们在函数 $s_i(\tau;U_-)$ 的所有非退化极小值这一假设下，验证\eqref{align:ch08-liu} 蕴涵\eqref{align:ch08-liu对偶}。（一般情形可由扰动论证化归为上述情形，逆命题的证明类似。）为确定起见，设 $\tau_+>0$，此时 $\tau_-<0$。由\eqref{align:ch08-liu} 可知，要么 $\dot s_i(\tau_+;U_-)<0$，要么 $\dot s_i(\tau_+;U_-)=0$ 且 $\ddot s_i(\tau_+;U_-)>0$。于是由\eqref{align:ch08-sdotpos}、\eqref{align:ch08-sdotzero} 与\eqref{align:ch08-sddot}，要么 $\lambda_i(U_+)<s$，要么 $\lambda_i(U_+)=s$ 且 $\mathrm{D}\lambda_i(U_+)R_i(U_+)>0$。在任一情形下，回顾\eqref{align:ch08-si0} 与\eqref{align:ch08-sdot}，可推得对接近 $0$ 的 $\tau<0$ 有 $s_i(\tau;U_+)<s$。若对某个 $\tau\in(\tau_-,0)$，\eqref{align:ch08-liu对偶} 被破坏，则 $s_i(\tau;U_+)-s$ 必在某个 $\tau_0\in(\tau_-,0)$ 处变号。令 $U_0=W_i(\tau_0;U_+)$。由于 $U_-$ 与 $U_0$ 都由速度为 $s$ 的激波连到 $U_+$，$U_-$ 与 $U_0$ 之间也可由速度为 $s$ 的激波相连，即对某个 $\tau_1\in(0,\tau_+)$ 有 $U_0=W_i(\tau_1;U_-)$，且 $s_i(\tau_1;U_-)=s$。于是\eqref{align:ch08-liu} 蕴涵 $\dot s_i(\tau_1;U_-)=0$ 且 $\ddot s_i(\tau_1;U_-)>0$，进而由\eqref{align:ch08-sdotzero}、\eqref{align:ch08-sddot} 有 $\lambda_i(U_0)=s$ 且 $\mathrm{D}\lambda_i(U_0)R_i(U_0)>0$。但此时\eqref{align:ch08-sdotzero}、\eqref{align:ch08-sddot} 再次给出 $\dot s_i(\tau_0;U_+)=0$ 且 $\ddot s_i(\tau_0;U_+)<0$，于是与假设相反，$s_i(\tau;U_+)-s$ 不能跨 $\tau_0$ 变号。

特别地，对 $\tau=0$ 应用\eqref{align:ch08-liu} 与\eqref{align:ch08-liu对偶} 并回顾\eqref{align:ch08-si0}，便得到\eqref{align:ch08-lax谱条件}。我们由此建立了

\begin{theorem}[Liu 准则蕴涵 Lax 准则]\label{theorem:ch08-liu蕴涵lax}
在\eqref{align:ch08-符号} 成立的范围内，任一满足 Liu $E$-条件的激波也满足 Lax $E$-条件。
\end{theorem}

当系统真非线性时，这两个准则至少在弱激波领域内重合：

\begin{theorem}[真非线性弱激波的 Liu 与 Lax 准则等价]\label{theorem:ch08-真非线性等价}
设 $i$-特征族真非线性且 $\lambda_i$ 是简单特征速度。则弱 $i$-激波满足 Liu $E$-条件，当且仅当它们满足 Lax $E$-条件。
\end{theorem}

\begin{proof}
由\cref{theorem:ch08-liu蕴涵lax}，Liu $E$-条件蕴涵 Lax $E$-条件。为证逆命题，设左侧状态 $U_-$ 由一个满足 Lax $E$-条件\eqref{align:ch08-lax谱条件} 的弱 $i$-激波（速度 $s$）在右侧连接到状态 $U_+=W_i(\tau_+;U_-)$。为确定起见，设 $\mathrm{D}\lambda_i(U_-)R_i(U_-)>0$（相反符号情形类似）。由\cref{theorem:ch08-弱激波lax}，$\tau_+<0$。由于激波为弱激波，由\cref{theorem:ch08-激波曲线}，在区间 $(\tau_+,0)$ 上 $\dot s_i(\tau;U_-)>0$。于是对 $\tau\in(\tau_+,0)$ 有 $s=s_i(\tau_+;U_-)<s_i(\tau;U_-)$，即 Liu $E$-条件成立。证明完毕。
\end{proof}

当系统非真非线性和/或激波非弱激波时，Liu $E$-条件比 Lax $E$-条件更严。这将通过下例说明。

首先考虑单个守恒律\eqref{align:ch07-标量守恒律}。激波曲线是 $u$ 轴，可直接以 $u$ 作为参数 $\xi$。激波速度由\eqref{align:ch08-标量s} 给出。于是显然，连接状态 $u_-$ 与 $u_+$ 的激波满足 Liu $E$-条件\eqref{align:ch08-liu}、\eqref{align:ch08-liu对偶}，当且仅当
\begin{align}
\frac{f(u_+)-f(u_0)}{u_+-u_0}\leqslant\frac{f(u_+)-f(u_-)}{u_+-u_-}\leqslant\frac{f(u_0)-f(u_-)}{u_0-u_-}
\label{align:ch08-oleinik}
\end{align}
对介于 $u_-$ 与 $u_+$ 之间的每个 $u_0$ 成立。这就是著名的 Oleinik $E$-条件。它很容易被记忆为一个几何命题：当 $u_-<u_+$（或 $u_->u_+$）时，将 $u_-$（左）连接到 $u_+$（右）的激波是容许的，当且仅当 $f$ 的图像以 $(u_-,f(u_-))$、$(u_+,f(u_+))$ 为端点的弧位于连接这两点的弦的上方（或下方）。令 $u_0$ 趋于 $u_-$ 与 $u_+$，可推得\eqref{align:ch08-oleinik} 蕴涵\eqref{align:ch08-标量lax}。当然，除非 $f$ 凸或凹，逆命题一般不成立。于是我们证明了：在单个守恒律中，当 $f$ 含有拐点时，Liu $E$-条件比 Lax $E$-条件更严。在真非线性情形，Liu 与 Lax $E$-条件等价。

现在转向等熵弹性系统\eqref{align:ch07-p系统}。激波曲线由\eqref{align:ch08-弹性轨迹} 确定，故可用 $u$ 代替 $\xi$ 作为参数。激波速度由\eqref{align:ch08-弹性s} 给出。因此，连接状态 $(u_-,v_-)$ 与 $(u_+,v_+)$ 的激波满足 Liu $E$-条件\eqref{align:ch08-liu}、\eqref{align:ch08-liu对偶}，当且仅当
\begin{align}
\frac{\sigma(u_0)-\sigma(u_-)}{u_0-u_-}\gtrless\frac{\sigma(u_+)-\sigma(u_-)}{u_+-u_-}\lessgtr\frac{\sigma(u_+)-\sigma(u_0)}{u_+-u_0}
\label{align:ch08-wendroff}
\end{align}
对介于 $u_-$ 与 $u_+$ 之间的所有 $u_0$ 成立，其中对 1-激波取“$\leqslant$”、对 2-激波取“$\geqslant$”。这称为 Wendroff $E$-条件。用几何语言可表述为：当 $s(u_+-u_-)<0$（或 $>0$）时，将 $(u_-,v_-)$（左）连接到 $(u_+,v_+)$（右）的激波是容许的，当且仅当 $\sigma$ 的图像以 $(u_-,\sigma(u_-))$、$(u_+,\sigma(u_+))$ 为端点的弧位于连接这两点的弦的下方（或上方）。显然，这与 Oleinik $E$-条件密切类比。令\eqref{align:ch08-wendroff} 中的 $u_0$ 趋于 $u_-$ 与 $u_+$，可推得 Wendroff $E$-条件蕴涵 Lax $E$-条件\eqref{align:ch08-压缩}。当 $\sigma$ 凸或凹时逆命题成立，否则不成立。因此，对系统\eqref{align:ch07-p系统}，当 $\sigma$ 含有拐点时，Liu $E$-条件比 Lax $E$-条件更严。在真非线性情形，Liu 与 Lax $E$-条件等价。

正如我们将看到的，Oleinik $E$-条件与 Wendroff $E$-条件自然地从其他容许性准则导出。在很大程度上，这些特殊 $E$-条件为一般 Liu $E$-条件的提出提供了动机。

\end{document}
